\documentclass[12pt]{article}
\usepackage[utf8]{inputenc}
\usepackage[T1]{fontenc}
\usepackage{hyperref}
\title{Samarokena Language Resources for K-12 Curriculum Development}
\author{Compiled by Claude Sonnet 5 (claude-sonnet-5) for Andr\'as Kornai}
\date{}
\begin{document}
\maketitle

\section{Introduction}
Samarokena (ISO 639-3 \texttt{tmj}), also called Samarkena, Karfasia or Tamaja, is a poorly documented Papuan language of Indonesian Papua. It is spoken in Karfasia, Maseb, Samarkena and Tamaya villages. Wikipedia cites about 400 speakers (1982). Its classification is unsettled.

\section{Primary Educational Resources}
None found.

\section{Linguistic and Reference Materials}
\begin{itemize}
\item Wikipedia, Samarokena language: \url{https://en.wikipedia.org/wiki/Samarokena_language}
\item Wikipedia, Kwerbic languages: \url{https://en.wikipedia.org/wiki/Kwerbic_languages}
\end{itemize}

\section{Religious and Community Resources}
Joshua Project reports that Bible translation has started. I could not confirm any published text. \url{https://joshuaproject.net/people_groups/14670}

\section{Digital and Audio Resources}
None found.

\section{Curriculum Development Context}
The search results contain no school materials, orthography or dictionary. Documentation would have to come before any curriculum work.
\end{document}